\documentclass[aspectratio=169]{beamer}

% Use a modern theme if available; keep a vivid fallback.
\IfFileExists{beamerthememetropolis.sty}{
  \usetheme[numbering=fraction,progressbar=foot]{metropolis}
}{
  \usetheme{Berlin}
  \usecolortheme{dolphin}
}
\setbeamertemplate{navigation symbols}{}

\usepackage[utf8]{inputenc}
\usepackage[T1]{fontenc}
\usepackage{lmodern}
\usepackage{booktabs}
\usepackage{tabularx}
\usepackage{array}
\usepackage{colortbl}
\usepackage{amsmath}
\usepackage{ragged2e}
\usepackage{xcolor}

% Vibrant custom palette
\definecolor{VividBlue}{RGB}{0,102,255}
\definecolor{VividOrange}{RGB}{255,111,0}
\definecolor{VividTeal}{RGB}{0,170,160}
\definecolor{DeepInk}{RGB}{20,28,44}

% Readability and visual consistency tweaks.
\setlength{\parskip}{0.25em}
\setlength{\parindent}{0pt}
\renewcommand{\arraystretch}{1.15}
\setlength{\tabcolsep}{6pt}
\setbeamertemplate{itemize item}{\small\textbullet}
\setbeamertemplate{itemize subitem}{\scriptsize\textbullet}
\setbeamercolor{normal text}{fg=DeepInk,bg=white}
\setbeamercolor{title}{fg=VividBlue}
\setbeamercolor{frametitle}{fg=white,bg=VividBlue}
\setbeamercolor{alerted text}{fg=VividOrange}
\setbeamercolor{structure}{fg=VividBlue}
\setbeamercolor{progress bar}{fg=VividOrange,bg=VividBlue!20}
\setbeamercolor{itemize item}{fg=VividOrange}
\setbeamercolor{itemize subitem}{fg=VividTeal}
\setbeamercolor{block title}{fg=white,bg=VividBlue}
\setbeamercolor{block body}{fg=DeepInk,bg=VividBlue!6}
\setbeamerfont{frametitle}{series=\bfseries,size=\large}
\setbeamerfont{title}{series=\bfseries}
\setbeamerfont{footline}{size=\scriptsize}

% Uncomment if you want speaker notes on a second screen:
% \usepackage{pgfpages}
% \setbeameroption{show notes on second screen=right}

\title{Salience-Guided Question Generation:\\Feature Ablations, Methodological Insights, and Issue Analysis}
\author{Madduru Sai Chandra Nikhil}
\date{April 2026}

\begin{document}

\begin{frame}
  \titlepage
\end{frame}

\begin{frame}{Agenda}
\begin{itemize}
  \item Problem \& Motivation
  \item Proposed Approach \& Features
  \item Experimental Setup \& Results
  \item Qualitative Analysis \& Insights
  \item Limitations \& Future Work
\end{itemize}
\end{frame}

\begin{frame}{Problem \& Motivation}
\small
\begin{itemize}
  \item Standard question generation is often fluent but poorly grounded -- questions may ignore key evidence in the passage.
  \item \textbf{SQuAD benchmark}: 100,000+ QA pairs, 3--5 questions per passage, answers are exact extractive spans.
  \item Challenge: How to steer the LLM toward the most answer-bearing sentences?
  \item Goal: Improve grounding and answer consistency using sentence-level salience guidance.
\end{itemize}

\scriptsize
\textit{Evaluated on 750 SQuAD passages.}
\end{frame}

\begin{frame}{Salience-Guided QG Pipeline}
\small
\begin{enumerate}
  \item Score every sentence using linguistic + surprisal features.
  \item Select top-k most salient (answer-bearing) sentences.
  \item Generate questions from only the selected salient evidence using Qwen2.5-7B-Instruct.
  \item Evaluate with overlap metrics + QA consistency.
\end{enumerate}

\begin{block}{Core Design Principle}
Soft grounding via salience guidance -- prioritize important sentences while keeping questions fluent and diverse.
\end{block}
\end{frame}

\begin{frame}{Salience Feature Engineering}
\small
\begin{tabularx}{\textwidth}{@{}X|X@{}}
\toprule
\textbf{Feature Family} & \textbf{Rationale} \\
\midrule
Linguistic (length, position, POS ratios, lexical density, discourse markers) & Interpretable, fast, strong structural priors for important sentences. \\
GPT-2 surprisal stats & Captures local uncertainty and information density. \\
BERT surprisal stats & Captures contextual uncertainty with bidirectional context. \\
\bottomrule
\end{tabularx}
\end{frame}

\begin{frame}{Linguistic Features (14 total)}
\small
\begin{tabularx}{\textwidth}{@{}lX@{}}
\toprule
\textbf{Category} & \textbf{Features} \\
\midrule
Surface & sentence\_length\_words, sentence\_position, sentence\_position\_norm, avg\_word\_length \\
Lexical & type\_token\_ratio, lexical\_density \\
POS ratios & noun\_ratio, verb\_ratio, adj\_ratio, pronoun\_ratio, noun\_verb\_ratio \\
Discourse & named\_entity\_density, causal\_marker\_ratio, contrast\_marker\_ratio \\
\bottomrule
\end{tabularx}
\end{frame}

\begin{frame}{Experimental Setup (750 Passages)}
\small
\textbf{Three salience conditions (same prompt, same LLM):}
\begin{itemize}
  \item \textbf{Full Salience}: 26 features (linguistic + surprisal)
  \item \textbf{Linguistic Only}: 14 linguistic features
  \item \textbf{Top-10}: Compact feature set
  \item \textbf{Baseline}: No salience guidance (random sentences)
\end{itemize}

\scriptsize Same evaluation pipeline for all conditions.
\end{frame}

\begin{frame}{Salience Classifier: Classification Metrics (2000 Passages)}
\small
\centering
\begin{tabular}{lccccc}
\toprule
\textbf{Configuration} & \textbf{Features} & \textbf{Accuracy} & \textbf{Precision} & \textbf{Recall} & \textbf{F1-Score} \\
\midrule
\textbf{Full (Ling + Surprisal)} & 26 & 0.693 & \textbf{0.786} & 0.695 & 0.738 \\
No Surprisal (Ling only) & 14 & 0.692 & 0.776 & \textbf{0.710} & \textbf{0.742} \\
Top-10 Compact & 10 & 0.674 & 0.765 & 0.687 & 0.724 \\
\bottomrule
\end{tabular}

\vspace{0.4em}
\textbf{Quick Summary:} Linguistic features are strong; surprisal mainly boosts precision.
\end{frame}

\begin{frame}{Feature Importance: What Drives Salience?}
\small
\centering
\textbf{Top-ranked features from the 2000-passage salience classifier}

\vspace{0.3em}
\includegraphics[width=0.92\linewidth,height=0.55\textheight,keepaspectratio]{../results/run_2000_passages/feature_importance.png}

\vspace{0.3em}
\textbf{Quick Summary:} The ranking shows that simple linguistic cues remain highly informative, while surprisal features add complementary signal for salience prediction.
\end{frame}

\begin{frame}{Question Generation Results (750 Passages)}
\centering
\textbf{Salience vs Baseline -- 5 Key Metrics}

\scriptsize Generated with Qwen2.5-7B-Instruct | Evaluated with deepset/roberta-base-squad2 (QA) \& roberta-base (BERTScore)

\vspace{0.2em}
\small
\begin{tabular}{l|ccc|ccc}
\toprule
& \multicolumn{3}{c|}{\textbf{Salience}} & \multicolumn{3}{c}{\textbf{Baseline}} \\
\textbf{Metric} & Full & NoS & Top10 & Full & NoS & Top10 \\
\midrule
ROUGE-L & 0.349 & 0.360 & 0.363 & 0.350 & 0.369 & \textbf{0.371} \\
QA-EM & 0.537 & 0.513 & 0.516 & 0.569 & \textbf{0.578} & 0.561 \\
QA-F1 & 0.775 & 0.764 & 0.772 & 0.767 & \textbf{0.778} & 0.765 \\
BERTScore & 0.905 & 0.908 & 0.908 & 0.906 & 0.909 & \textbf{0.910} \\
Answer-Recovery-F1 & 0.402 & 0.430 & 0.425 & 0.415 & \textbf{0.454} & 0.447 \\
\bottomrule
\end{tabular}

\vspace{0.3em}
\textbf{Key Observation:} Baseline wins on surface overlap. Salience wins on answer consistency.
\end{frame}

\begin{frame}{Salience vs. Baseline Metric Comparison}
\small
\begin{columns}[T,onlytextwidth]
\column{0.48\textwidth}
\begin{block}{Baseline Strength}
\begin{itemize}
  \item Better lexical overlap \& BERTScore
  \item Closer surface match to gold questions
\end{itemize}
\end{block}

\column{0.48\textwidth}
\begin{alertblock}{Salience Strength}
\begin{itemize}
  \item Better QA consistency \& answer recovery
  \item Stronger grounding to important evidence
\end{itemize}
\end{alertblock}
\end{columns}

\vspace{0.5em}
\centering\textbf{Takeaway: Grounding is not the same as gold overlap.}
\end{frame}

\begin{frame}{Why Questions Get Flagged}
\small
\textbf{Main reasons (from 2302 issues in qa\_generation\_issues.json)}

\begin{itemize}
  \item Low lexical overlap with gold question
  \item QA model answer does not match generated answer
  \item Generated question targets a different (but valid) fact
  \item QA model answer weakly aligned with gold answers
\end{itemize}

\textbf{Key Insight:} Many flagged questions are not low-quality -- they simply focus on different valid facts from the salient sentences.
\end{frame}

\begin{frame}{Data-Sourced Example 1: Better Semantic Alignment}
\small
\textbf{Warsaw\_Pact (example\_index 3) — Seeking Specificity}

\scriptsize \textbf{Gold (SQuAD question):} The UK was initially excluded from participating in which agreement?

\vspace{0.3em}
\begin{columns}[T,onlytextwidth]
\column{0.48\textwidth}
\begin{block}{Baseline}
Question: Why did the Soviets propose a new agreement? \\
Answer: to accept the participation of the USA \\
QA consistency F1: \textbf{0.67}
\end{block}

\column{0.48\textwidth}
\begin{alertblock}{Salience (Ours)}
Question: What did the Soviets declare about the NATO? \\
Answer: readiness to examine jointly with other parties \\
QA consistency F1: \textbf{0.94}
\end{alertblock}
\end{columns}

\vspace{0.3em}
\scriptsize \textit{Salient sentence}: ``The Soviets declared their readiness to examine jointly...''

\vspace{0.2em}
\textbf{Insight:} Salience grounds on more specific treaty language, improving QA consistency by 40\%.
\scriptsize Source: qa\_generation\_issues.json (example\_index 3)
\end{frame}

\begin{frame}{Data-Sourced Example 2: Grounding Recovers Details}
\small
\textbf{Warsaw\_Pact (example\_index 10) — Treaty Framework}

\scriptsize \textbf{Gold (SQuAD question):} How many countries formed the initial Warsaw Pact membership?

\vspace{0.3em}
\begin{columns}[T,onlytextwidth]
\column{0.48\textwidth}
\begin{block}{Baseline}
Question: What did Warsaw Pact members pledge? \\
Answer: mutual defense \\
QA consistency F1: \textbf{0.36}
\end{block}

\column{0.48\textwidth}
\begin{alertblock}{Salience (Ours)}
Question: What were treaty relations based on? \\
Answer: mutual non-intervention, sovereignty, independence \\
QA consistency F1: \textbf{0.70}
\end{alertblock}
\end{columns}

\vspace{0.3em}
\scriptsize \textit{Salient sentence}: ``Relations among treaty signatories were based on mutual non-intervention...''

\vspace{0.2em}
\textbf{Insight:} Salience captures nuanced treaty principles; QA consistency improves by 95\% via better passage grounding.
\scriptsize Source: qa\_generation\_issues.json (example\_index 10)
\end{frame}

\begin{frame}{Issue-Present Case: Both Methods Flagged}
\small
\textbf{Pub (example\_index 414)}

\scriptsize \textbf{Gold (SQuAD question):} In what historical period was a large portion of the population illiterate?

\vspace{0.25em}
\scriptsize \textit{Salient sentence}: ``For this reason there was often no reason to write the establishment's name on the sign...''

\vspace{0.3em}
\begin{columns}[T,onlytextwidth]
\column{0.48\textwidth}
\begin{block}{Baseline}
Question: Why were pictures useful on signs? \\
Answer: pictures on a sign were more useful than words \\
QA consistency F1: \textbf{0.0}
\end{block}

\column{0.48\textwidth}
\begin{alertblock}{Salience (Ours)}
Question: Why was writing the establishment's name on the sign uncommon? \\
Answer: pictures on a sign were more useful \\
QA consistency F1: \textbf{0.0}
\end{alertblock}
\end{columns}

\vspace{0.25em}
\textbf{Takeaway:} This is a true issue case: both methods ask reasonable questions, but they drift from the gold fact target, causing QA mismatch.
\scriptsize Source: qa\_generation\_issues.json (example\_index 414)
\end{frame}

\begin{frame}{Key Insights}
\small
\begin{enumerate}
  \item Full features give best consistency; compact features give better overlap.
  \item Linguistic features alone are surprisingly powerful.
  \item Grounding \(\neq\) gold overlap -- salience targets important content.
\end{enumerate}
\end{frame}

\begin{frame}{Future Work}
\small
\begin{block}{Next: Discourse-Level Salience with RST}
Current method is sentence-level only.  
Integrating Rhetorical Structure Theory (RST) will allow us to:
\begin{itemize}
  \item Identify nuclear (central) vs satellite (supporting) sentences
  \item Capture discourse relations (evidence, elaboration, contrast)
  \item Model multi-sentence evidence chains
\end{itemize}
\end{block}

Expected impact: significantly better salience ranking and consistency on complex questions.
\end{frame}

\begin{frame}{Conclusion \& Open Questions}
\small
\begin{itemize}
  \item Salience-guided QG is useful when grounding matters.
  \item Full features improve consistency; compact features improve overlap.
  \item Next step: optimize for the target objective, not one metric.
\end{itemize}
\end{frame}

\begin{frame}{Thank You / Q\&A}
\centering
\Large Thank you.

\vspace{1em}
\normalsize
Questions, feedback, and suggestions for the next iteration are very welcome.
\end{frame}

\end{document}
